\chapter{Gaussian elimination and higher state operations}
\label{ch:vol2-gaussian-elimination}

Let $p$ be expectation for a centred real Gaussian variable of
variance $\sigma>0$. Then $p(x)=0$ and $p(x^2)=\sigma$.
Eliminating $x$ therefore leaves information in products that its
mean does not contain. The connected contractions of
Chapter~\ref{ch:vol2-supported-compositions} exhibit the same effect
for fields. A reduction that preserves their operations must retain
these fluctuations together with the differential and the
normalization of the state.

In a finite polynomial model, an explicit homotopy identifies
each observable with its expectation. The failure to preserve
products becomes part of the comparison map. Formal interactions
change both the expectation and the differential of the variables
that remain. Successive integrations must preserve this changed
differential and the state on every product. We construct these
maps and compare them under changes of coordinates and of the
integration data.

The normal direction of a boundary comparison has an additional
datum: its endpoint condition. We first isolate that elementary
support question. The finite quantum construction that follows
does not require a choice of continuum propagator.

\section{A supported normal contraction}
\label{sec:vol2-normal-contraction}

Put $I=[0,\infty)$ and work over $\mathbb C$. The relative compactly
supported de Rham complex
\[
 C=\Omega_c^\bullet(I,\{0\})
\]
has $C^0=\{f\in C_c^\infty(I):f(0)=0\}$ and
$C^1=C_c^\infty(I)\,dt$, with $df=f'\,dt$.
Thus the relative condition constrains functions at the endpoint;
it imposes no restriction on the coefficient of a one-form there.
Choose $\rho\in C_c^\infty(I)$ with
$\int_0^\infty\rho(t)\,dt=1$.

Define maps to and from $\mathbb C[-1]$, placed in cohomological
degree one, by
\begin{equation}\label{eq:vol2-normal-maps}
 \begin{aligned}
 p(\beta\,dt)&=\int_0^\infty\beta(t)\,dt,& p(f)&=0,\\
 i(a)&=a\rho(t)\,dt,& H(f)&=0,\\
 H(\beta\,dt)(t)
   &=\int_0^t\left(\beta(s)-\rho(s)
                    \int_0^\infty\beta(r)\,dr\right)ds.
 \end{aligned}
\end{equation}
The subtraction in the primitive removes its constant tail.
It is required to keep the homotopy in the stated support class.

\begin{proposition}\label{prop:vol2-normal-retract}
The maps in \eqref{eq:vol2-normal-maps} are continuous and satisfy
\[
 pi=1,\qquad dH+Hd=1-ip,\qquad
 H^2=Hi=pH=0.
\]
The normal support of $H\beta$ lies in the compact interval spanned
by $0$, $\operatorname{supp}\beta$ and $\operatorname{supp}\rho$.
Tensoring this contraction with the full Dolbeault complex on an
open set $U\subset\mathbb C^d$ retains every tangential form degree.
\end{proposition}
\begin{proof}
The integrand defining $H\beta$ has total integral zero. Its
primitive vanishes at zero and again beyond the stated interval.
Differentiation gives $dH(\beta\,dt)=\beta\,dt-ip(\beta\,dt)$.
For $f\in C^0$, compact support and $f(0)=0$ give
\[
 p(df)=f(\infty)-f(0)=0,\qquad H(df)=f.
\]
These are the homotopy identities in the two degrees. The equation
$pi=1$ is the normalization of $\rho$; substituting $i(a)$ in $H$
gives $Hi=0$. The remaining side identities follow from degrees.

On functions or forms supported in a fixed compact interval,
integration is bounded by a finite smooth seminorm. The same
estimate for the primitive and its derivatives proves continuity
into the enlarged compact interval. These estimates give
continuity for the usual inductive-limit topology of compact
supports. On
$C\widehat\otimes_\pi\Omega^{0,\bullet}(U)$, the completed
projective tensor product, the maps extend by continuity.
The tensor differential is
\[
 d(\alpha\otimes\gamma)
 =d\alpha\otimes\gamma+(-1)^{|\alpha|}
                         \alpha\otimes\bar\partial\gamma.
\]
Its cross terms with $H\otimes1$ cancel. The resulting target is
$\Omega^{0,\bullet}(U)[-1]$, with differential $-\bar\partial$,
as required by the shift.
\end{proof}

Changing the endpoint condition changes the answer. In the
absolute complex $\Omega_c^\bullet(I)$, the primitive
\[
 H_{\mathrm{abs}}(\beta\,dt)(t)
   =-\int_t^\infty\beta(s)\,ds
\]
is allowed even when it is nonzero at $t=0$. It is compactly
supported and satisfies $dH_{\mathrm{abs}}=1$ on one-forms and
$H_{\mathrm{abs}}d=1$ on functions. The absolute complex is
acyclic. In a field theory, the polarization must determine which
normal complex occurs before either contraction is used.

\section{The finite Gaussian Ward complex}
\label{sec:vol2-gaussian-ward}

Let $R_0$ be a nonzero commutative $\mathbb Q$-algebra and put
$R=R_0[\hbar]$, where $\hbar$ has degree zero. Fix a symmetric
invertible matrix $A\in\operatorname{Mat}_m(R_0)$ and write
$C=A^{-1}$. Introduce even coordinates $q_1,\ldots,q_m$ of
degree zero and odd coordinates $\eta_1,\ldots,\eta_m$ of
degree $-1$. Our observable algebra is
\[
 \mathcal G_A=R[q_1,\ldots,q_m]\otimes
                         \Lambda(\eta_1,\ldots,\eta_m).
\]
All odd derivatives act on the left. Tensor products and
reorderings use the cohomological Koszul sign.

The quadratic potential $S_0(q)=\tfrac12q^{\mathrm t}Aq$
defines the \emph{Ward differential}
\begin{equation}\label{eq:vol2-gaussian-ward}
 D_A=\sum_i\bigl((Aq)_i-\hbar\partial_{q_i}\bigr)
                                             \partial_{\eta_i}.
\end{equation}
This is the finite-dimensional homological form of Gaussian
integration by parts; see Gwilliam--Johnson-Freyd
\cite{GwilliamJohnsonFreyd2019}*{\S\S1--3}.
The even operators $(Aq)_i-\hbar\partial_{q_i}$ commute, because
$A$ is symmetric. The odd derivatives anticommute, so $D_A^2=0$.
In one variable, with $A=(a)$,
\begin{equation}\label{eq:vol2-ward-square}
 D_A\eta=ax,\qquad
 D_A(\eta x)=ax^2-\hbar.
\end{equation}
The second equality already records the fluctuation that survives
elimination.

We will construct a \emph{contraction} of this complex onto $R$:
cochain maps $i:R\to\mathcal G_A$, $p:\mathcal G_A\to R$ and
a degree-minus-one map $H$ satisfying
\[
 pi=1,\qquad D_AH+HD_A=1-ip,
 \qquad H^2=Hi=pH=0.
\]
The last three equations are the side conditions of a strong
deformation retract. The projection $p$, normalized by $p(1)=1$,
will be the formal Gaussian state.

\subsection{Heat conjugation and the Euler homotopy}

Let
\[
 d_0=\sum_i(Aq)_i\partial_{\eta_i},\qquad
 T_A=\exp\!\left(\frac{\hbar}{2}
                 \sum_{i,j}C_{ij}\partial_{q_i}\partial_{q_j}\right).
\]
The heat series and its inverse terminate on every polynomial.
The commutator of its exponent with $q_i$ is
$\hbar\sum_jC_{ij}\partial_{q_j}$, and the next commutator is zero.
Consequently
\begin{equation}\label{eq:vol2-ward-conjugation}
 T_A^{-1}q_iT_A=q_i-\hbar\sum_jC_{ij}\partial_{q_j},
 \qquad D_A=T_A^{-1}d_0T_A.
\end{equation}

Let $N$ count total polynomial degree in the $q_i$ plus exterior
degree in the $\eta_i$, and put
\[
 R_A=\sum_{i,j}\eta_iC_{ij}\partial_{q_j},\qquad
 h_0=N^{-1}R_A\quad(N>0),\qquad h_0(1)=0.
\]
The notation $N^{-1}$ divides a homogeneous monomial by its
positive integer degree. These divisions are valid over $R_0$.
Let $\epsilon$ set every $q_i,\eta_i$ equal to zero, and let
$i$ include constants.

\begin{lemma}\label{lem:vol2-euler-contraction}
The maps $(\epsilon,i,h_0)$ contract $(\mathcal G_A,d_0)$ onto $R$.
\end{lemma}
\begin{proof}
The anticommutator of the odd derivations $d_0$ and $R_A$ is the
Euler operator:
\[
 d_0R_A+R_Ad_0
   =\sum_iq_i\partial_{q_i}
      +\sum_i\eta_i\partial_{\eta_i}=N.
\]
Indeed, its value on each $q_i$ and $\eta_i$ is that generator;
an even derivation is determined by these values. Both odd
operators preserve $N$. Division by $N$ proves
$d_0h_0+h_0d_0=1$ in positive total degree, while both sides of
$1-i\epsilon$ vanish on constants. The constant-coefficient even
derivatives in $R_A$ commute and its odd coefficients
anticommute, giving $R_A^2=0$. Hence $h_0^2=0$.
The equations $h_0i=\epsilon h_0=0$ follow from the same formulas.
\end{proof}

\begin{theorem}\label{thm:vol2-gaussian-retract}
The maps
\begin{equation}\label{eq:vol2-gaussian-retract}
 p_A=\epsilon T_A,\qquad i(a)=a,\qquad
 H_A=T_A^{-1}h_0T_A
\end{equation}
give a contraction of $(\mathcal G_A,D_A)$ onto $R$.
On even polynomials, $p_A$ is the Wick functional of covariance
$\hbar A^{-1}$; it kills every term of negative cohomological degree.
All maps are defined over $R_0[\hbar]$.
\end{theorem}
\begin{proof}
Conjugate every identity of Lemma~\ref{lem:vol2-euler-contraction}
by $T_A$. Equation~\eqref{eq:vol2-ward-conjugation} gives the
new differential, and $T_Ai=i$ keeps the inclusion fixed.
Thus the homotopy and all side identities follow on every
polynomial, without a convergence argument.

To identify the state, introduce independent degree-zero formal
variables $J_1,\ldots,J_m$. Applying $\epsilon T_A$
coefficientwise gives
\begin{equation}\label{eq:vol2-wick-generating}
 p_A(e^{J^{\mathrm t}q})
       =\exp\!\left(\frac{\hbar}{2}J^{\mathrm t}CJ\right).
\end{equation}
Differentiation in the $J_i$ pairs all input coordinates, with
$\hbar C_{ij}$ for each pair. This proves Wick's formula and
determines $p_A$ on every polynomial. Heat operators preserve
exterior degree, so evaluation kills negative cohomological degrees.
\end{proof}

For $A=(a)$, the retract gives the concrete primitive
\begin{equation}\label{eq:vol2-quadratic-primitive}
 p_A(x^2)=\frac{\hbar}{a},\qquad
 H_A(x^2)=\frac{\eta x}{a},\qquad
 D_A\left(\frac{\eta x}{a}\right)=x^2-\frac{\hbar}{a}.
\end{equation}
At a positive real specialization $A>0$, $\hbar>0$,
equation~\eqref{eq:vol2-wick-generating} is also the generating
function of the normalized convergent Gaussian integral.
Positivity is needed for that real integral, and is not a
hypothesis of the algebraic contraction.

\subsection{Ordinary and transported products}

The ordinary product on $\mathcal G_A$ is not a cochain operation:
\eqref{eq:vol2-ward-square} gives
\[
 D_A(\eta x)-(D_A\eta)x=-\hbar.
\]
One can transport a product along the heat isomorphism. It is
useful to do this explicitly, since the transported square has a
different expectation.

\begin{proposition}\label{prop:vol2-transported-product}
In one variable define
\begin{equation}\label{eq:vol2-transported-product}
 f\diamond g=T_A^{-1}\bigl((T_Af)(T_Ag)\bigr)
 =\sum_{r\geq0}\frac{(-\hbar/a)^r}{r!}
                  (\partial_x^rf)(\partial_x^rg).
\end{equation}
This is a unital graded-commutative associative product for which
$D_A$ is a derivation. The map $p_A$ is an augmentation for
$\diamond$ and a quasi-isomorphism. In particular,
\[
 x\diamond x=x^2-\hbar/a,\qquad
 p_A(x\diamond x)=0,\qquad p_A(x^2)=\hbar/a.
\]
\end{proposition}
\begin{proof}
The differential $d_0=ax\partial_\eta$ is a derivation for the
ordinary product. Transporting that differential and product by
$T_A^{-1}$ proves the algebra and differential identities.
Evaluation at zero is an augmentation of the classical Koszul
complex, so its transported map is $p_A$.

To obtain the sum in \eqref{eq:vol2-transported-product}, apply
the heat operators to two independent variables $u,v$ before
setting $u=v=x$. The second derivative of a product is induced
by $(\partial_u+\partial_v)^2$. Its two individual square terms
cancel those of $T_A$ on the factors, leaving
$\exp(-\hbar\partial_u\partial_v/a)$. All these series terminate
on polynomials. The displayed square follows at once, and the
quasi-isomorphism follows from Theorem~\ref{thm:vol2-gaussian-retract}.
\end{proof}

For the ordinary observables, the variance in
\eqref{eq:vol2-quadratic-primitive} still has to be retained.
Its higher analogues give a comparison of products without
changing which polynomial is being observed.

\section{Cumulants as operation maps}
\label{sec:vol2-state-cumulants}

Let $B$ be a unital graded-commutative algebra over a commutative
$\mathbb Q$-algebra. Write
\[
 S^c(B)=\bigoplus_{n\geq0}\operatorname{Sym}^n(B)
\]
for the coaugmented conilpotent symmetric coalgebra of its
underlying graded module, with the unshuffle coproduct.
A coaugmentation-preserving coalgebra map out of $S^c(B)$ is
determined by its \emph{Taylor coefficients}: its composites
from each symmetric length to the length-one component of the
target. The scalar in length zero and the algebra unit
$1_B\in B$, which lies in length one, are distinct.

The products of $B$ determine a coalgebra automorphism $\tau_B$
whose $n$th Taylor coefficient is
\[
 (\tau_B)_n(b_1\odot\cdots\odot b_n)=b_1\cdots b_n,
 \qquad n\geq1.
\]
On an input of length $n$, it sums over partitions of
$\{1,\ldots,n\}$: multiply within each block and take the
symmetric tensor of the block products. Within a block the
indices are increasing; blocks are ordered by their least
indices. Reordering homogeneous inputs into this order gives
the Koszul sign. The term with singleton blocks is the identity,
and every other term lowers length. Thus the inverse is finite
on each input.

\begin{lemma}\label{lem:vol2-cumulant-inverse}
The $n$th Taylor coefficient of $\tau_B^{-1}$ is
\[
 (\tau_B^{-1})_n=(-1)^{n-1}(n-1)!\,\mu_n,
\]
where $\mu_n$ is multiplication of $n$ elements.
\end{lemma}
\begin{proof}
Corestrict the proposed composite $\tau_B^{-1}\tau_B$ to length
one. All partitions into $k$ blocks have the same coefficient
$(-1)^{k-1}(k-1)!$ after their products are multiplied.
The resulting scalar is
\[
 \sum_{k=1}^n
 \left\{\begin{matrix}n\\k\end{matrix}\right\}
                       (-1)^{k-1}(k-1)! =
 \begin{cases}1,&n=1,\\0,&n>1,\end{cases}
\]
where the braces count set partitions into $k$ nonempty blocks.
This identity is the coefficient identity of
$\log(e^z)=z$. The same regrouping supplies the Koszul signs
of the graded products. A coalgebra map is determined by these
corestrictions, proving the claimed inverse.
\end{proof}

\subsection{The deviations of a differential}

Suppose $D:B\to B$ has degree one, $D^2=0$ and $D1_B=0$.
It need not be a derivation. Let $\widetilde D$ be the unary
coderivation on $S^c(B)$: it applies $D$ to one factor at a
time, with the tensor signs. Put
\begin{equation}\label{eq:vol2-deviation-coderivation}
 \mathscr D_B=\tau_B^{-1}\widetilde D\tau_B.
\end{equation}
This is again a degree-one square-zero coderivation. Such a
coderivation with zero constant Taylor coefficient is an
$L_\infty[1]$ structure; its operations are graded symmetric
and have degree one.

\begin{proposition}\label{prop:vol2-multiplicative-deviations}
The first Taylor coefficients of \eqref{eq:vol2-deviation-coderivation}
are
\[
 q_1(b)=Db,\qquad
 q_2(b,c)=D(bc)-(Db)c-(-1)^{|b|}b(Dc).
\]
If $D$ has differential-operator order at most two, all Taylor
coefficients of arity at least three vanish.
\end{proposition}
\begin{proof}
In length two, the product term in $\tau_B$ followed by $D$
contributes $D(bc)$. The length-two part of $\tau_B^{-1}$
subtracts the two unary terms, giving $q_2$.
More generally, the Taylor coefficient is the iterated graded
commutator of $D$ with the multiplication operators of the
inputs, evaluated at $1_B$. By definition, an operator of order
at most two has zero commutators with three or more
multiplication operators. The square-zero identities for the
remaining operations follow from conjugation in
\eqref{eq:vol2-deviation-coderivation}.
\end{proof}

For $D_A$, the binary deviation is divisible by $\hbar$.
Equation~\eqref{eq:vol2-ward-square} says
$q_2(\eta,x)=-\hbar$ in one variable. This second-order
operation is part of the Ward complex.

\subsection{The higher components of a normalized map}

Let $P:(B,D_B)\to(C,D_C)$ be a degree-zero cochain map between
unital graded-commutative algebras, with $P1_B=1_C$.
The linear map need not preserve products. Define
\begin{equation}\label{eq:vol2-cumulant-coalgebra}
 \mathscr K(P)=\tau_C^{-1}S^c(P)\tau_B.
\end{equation}
Its $n$th Taylor coefficient is the \emph{cumulant}
\begin{equation}\label{eq:vol2-state-cumulant}
 \kappa_n^P(b_1,\ldots,b_n)=
 \sum_{\pi\in\Pi_n}(-1)^{|\pi|-1}(|\pi|-1)!
 \epsilon(\pi;b)
 \prod_{E\in\pi}P\!\left(\prod_{j\in E}b_j\right).
\end{equation}
Here $\Pi_n$ is the set of partitions into nonempty blocks and
$\epsilon(\pi;b)$ is the sign of the grouping specified above.
In particular,
\[
 \kappa_1^P=P,\qquad
 \kappa_2^P(b,c)=P(bc)-P(b)P(c).
\]
This coalgebra construction and its compatibility with
Koszul brackets are described by Bandiera
\cite{Bandiera2020}*{\S1.2, Lemma 1.16 and Proposition 1.21}.

\begin{theorem}\label{thm:vol2-cumulant-comparison}
The map $\mathscr K(P)$ is an $L_\infty[1]$ morphism between
the multiplicative deviations of $D_B$ and $D_C$:
\[
 \mathscr D_C\mathscr K(P)=\mathscr K(P)\mathscr D_B.
\]
If $P$ is a quasi-isomorphism, this is an $L_\infty[1]$
quasi-isomorphism. Every cumulant of arity greater than one
vanishes on an algebra-unit input. For a further normalized
cochain map $Q:C\to E$,
\begin{equation}\label{eq:vol2-cumulant-sewing}
 \mathscr K(QP)=\mathscr K(Q)\mathscr K(P).
\end{equation}
\end{theorem}
\begin{proof}
Inserting \eqref{eq:vol2-cumulant-coalgebra} and cancelling
$\tau_C\tau_C^{-1}$ reduces the morphism equation to
$\widetilde D_CS^c(P)=S^c(P)\widetilde D_B$.
This is precisely the cochain equation for $P$, applied to
each tensor factor. The linear Taylor coefficient is $P$;
being a quasi-isomorphism in that coefficient is the definition
of an $L_\infty$ quasi-isomorphism.

The unit assertion can be seen by adding independent marked
variables of the degrees that make a sum $F$ of marked inputs
have total degree zero. The cumulants are the coefficients of
$\log P(e^F)$. Adding a scalar mark $z1_B$ changes this expression
to $z+\log P(e^F)$ because $P1_B=1_C$.
All mixed coefficients involving $z$ therefore vanish.
This is a finite coefficient identity in any specified arity.

Finally the two adjacent multiplication automorphisms in
$\mathscr K(Q)\mathscr K(P)$ cancel. The remaining composite
is \eqref{eq:vol2-cumulant-coalgebra} for $QP$, proving
\eqref{eq:vol2-cumulant-sewing} in every arity.
\end{proof}

At arity two, \eqref{eq:vol2-cumulant-sewing} is the law of
total covariance:
\begin{equation}\label{eq:vol2-total-covariance}
 \kappa_2^{QP}(b,c)
   =Q\bigl(\kappa_2^P(b,c)\bigr)
                       +\kappa_2^Q(Pb,Pc).
\end{equation}
The second term measures fluctuations of the variables left
after the first elimination.

\subsection{A Gaussian state has higher operations on composites}

Take one variable of stiffness $a$ and write $p=p_A$.
For even polynomials $f,g_2,\ldots,g_n$, the morphism equation
of Theorem~\ref{thm:vol2-cumulant-comparison} on
$\eta f,g_2,\ldots,g_n$ gives, for $n\geq2$,
\begin{equation}\label{eq:vol2-higher-ward}
 \begin{split}
 \kappa_n^p(axf-\hbar f',g_2,\ldots,g_n)
  =\hbar\sum_{j=2}^n
    \kappa_{n-1}^p(fg_j',g_2,\ldots,\widehat g_j,\ldots,g_n).
 \end{split}
\end{equation}
Indeed the only binary input terms are
$q_2(\eta f,g_j)=-\hbar fg_j'$, and the target differential
and its deviations are zero. For $n=2$, $f=1$, $g_2=x$,
the equation forces $a\kappa_2^p(x,x)=\hbar$.
The variance is required by the cochain equation itself.

Although the higher cumulants of the linear variable vanish
above arity two, those of its square do not.
Wick's formula gives the formal generating function
\[
 p(e^{tx^2})=(1-2t\hbar/a)^{-1/2}.
\]
Taking its logarithm and comparing coefficients of $t^n/n!$
proves
\begin{equation}\label{eq:vol2-square-cumulants}
 \kappa_n^p(x^2,\ldots,x^2)
   =2^{n-1}(n-1)!\left(\frac{\hbar}{a}\right)^n,
 \qquad n\geq1.
\end{equation}
Thus the state on the full polynomial algebra has nonzero
components in arbitrarily high arity. Its Gaussian character
does not truncate the operations of composite observables.

\section{Formal interactions and conditional elimination}
\label{sec:vol2-interacting-elimination}

Introduce a degree-zero coupling parameter $g$. We work in
\[
 \bigl(R_0[\hbar,q]\otimes\Lambda\eta\bigr)[[g]].
\]
This means a polynomial in the even and odd coordinates at
each coupling order. Polynomial degree may grow with that order;
there is no completion in the field coordinates themselves.
The positive $g$-adic filtration makes the following perturbation
construction well defined.

\subsection{A filtered perturbation identity}

\begin{lemma}\label{lem:vol2-filtered-perturbation}
Let $(i,p,h)$ be a contraction of a complex $(C,d)$, complete
and separated for a descending filtration, onto $(B,d_B)$, with
\[
 dh+hd=1-ip,\qquad pi=1,\qquad h^2=hi=ph=0.
\]
Assume all these maps preserve filtration. Let $\delta$ have
degree one, strictly raise filtration, and satisfy
$(d+\delta)^2=0$. Put
\[
 F=(1+\delta h)^{-1}\delta=\delta(1+h\delta)^{-1}.
\]
Then a contraction for $d+\delta$ is given by
\begin{equation}\label{eq:vol2-perturbation-formulas}
 \begin{aligned}
 d_B'&=d_B+pFi,&
 i'&=i-hFi=(1+h\delta)^{-1}i,\\
 p'&=p-pFh=p(1+\delta h)^{-1},&
 h'&=h-hFh=h(1+\delta h)^{-1}.
 \end{aligned}
\end{equation}
In particular, $(d_B')^2=0$, $p'i'=1$ and
\[
 (d+\delta)h'+h'(d+\delta)=1-i'p',\qquad
 (h')^2=h'i'=p'h'=0.
\]
\end{lemma}
\begin{proof}
All inverses are geometric series. At any fixed filtration
precision only finitely many perturbation factors occur.
On endomorphisms write
$\partial T=dT-(-1)^{|T|}Td$.
The hypotheses give $\partial\delta=-\delta^2$ and
$\partial h=1-ip$. Apply $\partial$ to
$F=\delta-\delta hF$. The graded product rule gives
\[
 (1+\delta h)\partial F=-\delta ipF,
 \qquad\text{hence}\qquad
 dF+Fd=-FipF.
\]

It follows that
\[
 (d_B')^2=p(dF+Fd+FipF)i=0.
\]
Substitution in the inclusion equation gives
\[
 (d+\delta)i'-i'd_B'=h(dF+Fd+FipF)i=0.
\]
Using both $F=\delta-\delta hF$ and
$F=\delta-Fh\delta$, the homotopy equation expands to
\begin{align*}
 (d+\delta)h'+h'(d+\delta)
 &=1-ip+ipFh+hFip+h(dF+Fd)h\\
 &=1-(i-hFi)(p-pFh).
\end{align*}
The original side conditions give $p'i'=1$ and
$(h')^2=h'i'=p'h'=0$ directly. Commuting the homotopy
equation with $d+\delta$ and using the inclusion equation
gives $i'(d_B'p'-p'(d+\delta))=0$. Apply $p'$ to obtain
the projection equation. This proves all the contraction
identities.
\end{proof}

\subsection{An interacting state over the integral quantum ring}

Let
\[
 S(q)=\tfrac12q^{\mathrm t}Aq+V(q),\qquad
 V\in gR_0[\hbar,q][[g]].
\]
The potential is independent of the odd variables. Its Ward
differential is
\begin{equation}\label{eq:vol2-potential-ward}
 D_S=\sum_i(S_i-\hbar\partial_{q_i})\partial_{\eta_i},
 \qquad S_i=\partial_{q_i}S.
\end{equation}
The operators in parentheses commute, since
\[
 [S_i-\hbar\partial_{q_i},S_j-\hbar\partial_{q_j}]
     =\hbar(\partial_{q_j}S_i-\partial_{q_i}S_j)=0.
\]
Thus $D_S^2=0$. The perturbation
$\delta_V=\sum_iV_i\partial_{\eta_i}$ has positive coupling
order and kills constants.

\begin{theorem}\label{thm:vol2-formal-state-retract}
The Gaussian maps of Theorem~\ref{thm:vol2-gaussian-retract}
give an interacting contraction with
\begin{equation}\label{eq:vol2-formal-state-retract}
 p_V=p_A(1+\delta_VH_A)^{-1},\qquad
 i_V=i,\qquad
 H_V=H_A(1+\delta_VH_A)^{-1}.
\end{equation}
Its target is $R_0[\hbar][[g]]$ with zero differential.
All coefficients are polynomial in $\hbar$ at each coupling
order. The normalized exponential formula
\begin{equation}\label{eq:vol2-interacting-expectation}
 p_V(f)=\frac{p_A(fe^{-V/\hbar})}{p_A(e^{-V/\hbar})}
\end{equation}
agrees with this projection, with negative cohomological degrees
sent to zero. The quotient is first interpreted in
$R_0[\hbar,\hbar^{-1},q][[g]]$; its normalized value lies in
the integral quantum ring.
\end{theorem}
\begin{proof}
Apply Lemma~\ref{lem:vol2-filtered-perturbation} to the
Gaussian contraction. Since $\delta_Vi=0$, the inclusion and
target differential remain unchanged. The other formulas are
\eqref{eq:vol2-formal-state-retract}. Each coupling coefficient
contains finitely many operations on polynomials. None of these
operations divides by $\hbar$.

In the Laurent coefficient ring, the denominator of
\eqref{eq:vol2-interacting-expectation} is $1+O(g)$ and
therefore a unit. Gaussian integration by parts says
\[
 p_A\bigl((Aq)_if-\hbar\partial_{q_i}f\bigr)=0.
\]
Apply this identity to $fe^{-V/\hbar}$ at each coupling order.
The normalized quotient $L$ then satisfies
\[
 L\bigl(S_if-\hbar\partial_{q_i}f\bigr)=0,
 \qquad L(1)=1.
\]
These equations give the cochain identity on inputs with
one odd generator. With two or more odd generators,
the image under $D_S$ still has negative cohomological degree
and is killed by $L$.

The normalized cochain functional is unique. For an even
polynomial $f$, the free contraction gives
$f-ip_Af=D_AH_Af$. Hence any such functional satisfies
\[
 L(f)=p_Af-L(\delta_VH_Af).
\]
Solving recursively in coupling order gives
$L=p_A(1+\delta_VH_A)^{-1}$. This identifies the quotient
with the already integral projection.
\end{proof}

For the formal cubic potential $S=ax^2/2+gx^3/6$, the first
moments are
\begin{equation}\label{eq:vol2-cubic-state-moments}
 \begin{aligned}
 p_V(x)&=-\frac{g\hbar}{2a^2}
             -\frac{5g^3\hbar^2}{8a^5}+O(g^5),\\
 p_V(x^2)&=\frac{\hbar}{a}
             +\frac{5g^2\hbar^2}{4a^4}+O(g^4).
 \end{aligned}
\end{equation}
These follow by expanding \eqref{eq:vol2-interacting-expectation}
and using $p_A(x^{2n})=(2n-1)!!(\hbar/a)^n$.
They obey $a p_V(x)+(g/2)p_V(x^2)=0$, the Ward equation on
the odd generator. All series in this example are formal in $g$.

\subsection{The differential of the retained variables}

Split the even variables into $x=(x_1,\ldots,x_r)$ and
$y=(y_1,\ldots,y_s)$, with respective odd variables
$\eta_x,\eta_y$. Fix symmetric invertible matrices $A_x,A_y$
over $R_0$, and put
\begin{equation}\label{eq:vol2-conditional-potential}
 S(x,y)=\tfrac12x^{\mathrm t}A_xx
             +\tfrac12y^{\mathrm t}A_yy+V(x,y),
 \qquad V\in gR_0[\hbar,x,y][[g]].
\end{equation}
Write $P_x^0$ for the free Wick expectation in $x$. It kills
every term containing an $\eta_x$ and preserves the
$y,\eta_y$ variables. Define
\begin{equation}\label{eq:vol2-conditional-expectation}
 \begin{aligned}
 Z_x(y)&=P_x^0(e^{-V/\hbar}),&
 P_x(f)&=Z_x^{-1}P_x^0(fe^{-V/\hbar}),\\
 S_{\mathrm{eff}}(y)&=\tfrac12y^{\mathrm t}A_yy
                                      -\hbar\log Z_x(y).
 \end{aligned}
\end{equation}
These exponential and logarithmic expressions initially have
Laurent quantum coefficients and are complete in $g$.
The odd-variable rule is the same as for $P_x^0$.

For any potential $W$ in a tuple of even variables $q$, write
\[
 \mathcal B_W=
       \bigl(R_0[\hbar,q]\otimes\Lambda\eta_q\bigr)[[g]]
\]
with differential \eqref{eq:vol2-potential-ward}.
Its binary multiplicative deviation $q_2^W$ is divisible by
$\hbar$; denote the quotient by $\overline q_2^W$.
The operations $(D_W,\overline q_2^W)$ form an
$L_\infty[1]$ structure. Indeed the original square-zero
coderivation has only unary and binary components. Its
differential identity for the binary operation has one
factor of $\hbar$, and its Jacobi identity has two.
Multiplication by $\hbar$ is injective in these polynomial
coefficient modules, even if $R_0$ has zero divisors, so both
identities can be divided.

\begin{theorem}[Integral conditional elimination]
\label{thm:vol2-integral-conditional}
The effective potential belongs to $R_0[\hbar,y][[g]]$.
The normalized map in \eqref{eq:vol2-conditional-expectation}
is a unital cochain quasi-isomorphism
\[
 P_x:(\mathcal B_S,D_S)\longrightarrow
                       (\mathcal B_{S_{\mathrm{eff}}},D_{S_{\mathrm{eff}}})
\]
over $R_0[\hbar][[g]]$. It is the projection in the
filtered perturbation of the free Gaussian contraction.
For every $n\geq1$, its cumulant is divisible by
$\hbar^{n-1}$. The integral maps
\begin{equation}\label{eq:vol2-divided-cumulants}
 \Phi_n^{P_x}=\hbar^{1-n}\kappa_n^{P_x}
\end{equation}
give an $L_\infty[1]$ quasi-isomorphism for the divided
operations $(D_S,\overline q_2^S)$ and
$(D_{S_{\mathrm{eff}}},\overline q_2^{S_{\mathrm{eff}}})$.
They are continuous in $g$, polynomial in $\hbar$ at each
coupling order, and vanish on a unit input when $n>1$.
\end{theorem}
\begin{proof}
First work in Laurent quantum coefficients. The hidden Ward
equations are
\[
 P_x\bigl(S_{x_i}f-\hbar\partial_{x_i}f\bigr)=0.
\]
Differentiating the normalized quotient in a retained variable
gives
\begin{align}
 \partial_{y_j}P_x(f)
 &=P_x(\partial_{y_j}f)
       -\hbar^{-1}\bigl(P_x(fV_{y_j})-P_x(f)P_x(V_{y_j})\bigr),
       \label{eq:vol2-retained-covariance}\\
 (S_{\mathrm{eff}})_{y_j}
 &=(A_yy)_j+P_x(V_{y_j}).\notag
\end{align}
Combining these equations yields
\begin{equation}\label{eq:vol2-retained-ward}
 P_x(S_{y_j}f-\hbar\partial_{y_j}f)
   =(S_{\mathrm{eff}})_{y_j}P_x(f)
                           -\hbar\partial_{y_j}P_x(f).
\end{equation}
The hidden identity deals with the last eliminated odd
generator; \eqref{eq:vol2-retained-ward} deals with every
retained odd generator, with its left-exterior sign.
Terms still containing an eliminated odd generator vanish.
This proves the cochain equation in every exterior degree.

Let $(i_0,P_x^0,H_0)$ be the Gaussian contraction in $x$
tensored with the retained free Ward complex. Apply
Lemma~\ref{lem:vol2-filtered-perturbation} with
$\delta=\sum V_i\partial_{\eta_i}$. Its projection is
$P_x^0(1+\delta H_0)^{-1}$. Each $H_0$ adds one hidden
odd generator. A term surviving the final $P_x^0$ must
start without a hidden odd generator, and each subsequent
$\delta$ must remove the hidden generator just introduced.
Thus this projection acts on retained coefficients by a
normalized functional for the hidden Ward differential.
The uniqueness recursion used in
Theorem~\ref{thm:vol2-formal-state-retract}, now with $y$
as coefficients, identifies it with $P_x$.

In the transferred differential, one more $\delta$ occurs
than $H_0$. A surviving correction therefore removes exactly
one retained odd generator and otherwise only hidden ones.
It has the form $\sum_jF_j(y)\partial_{\eta_{y_j}}$.
Evaluating the cochain equation on $\eta_{y_j}$ identifies
$F_j=\partial_{y_j}(S_{\mathrm{eff}}-\tfrac12y^{\mathrm t}A_yy)$.
The retained second-order term is unchanged. The filtered
lemma now gives the asserted contraction and quasi-isomorphism.
Its projection, inclusion and homotopy are all integral in
$\hbar$.

It remains to prove integrality of the full effective potential
and the stated powers in its cumulants. Expand Wick
expectation with independent marked sources for the observables.
The logarithm selects connected graphs. A graph with $n$
marked vertices, $v$ interaction vertices and $E$ edges has
\[
 E\geq n+v-1.
\]
Every edge contributes a factor $\hbar A_x^{-1}$ and each
interaction vertex contributes at worst $\hbar^{-1}$.
The resulting power is at least $E-v\geq n-1$.
At coupling order $g^k$, the positive coupling orders of
the vertices sum to $k$, so $v\leq k$. For fixed polynomial
inputs there are only finitely many graphs at that order.
Retained variables are coefficients in this count, including
their exterior signs. An eliminated odd input gives zero.

Vacuum components cancel from normalized cumulants. For
$\log Z_x$, a connected vacuum graph satisfies $E\geq v-1$;
the factor $-\hbar$ in the effective potential removes its
one possible negative quantum power. An interaction independent
of $x$ contributes its original potential term, so this argument
also includes zero-valent vertices. It proves both remaining
integrality assertions.

Finally apply Theorem~\ref{thm:vol2-cumulant-comparison} to
$P_x$. In its arity-$n$ morphism equation, a unary term has
the factor $\hbar^{n-1}$. A source binary term has
$\hbar\hbar^{n-2}$, while a target binary term with input
arities $r+s=n$ has
$\hbar\hbar^{r-1}\hbar^{s-1}$. Dividing these common
factors gives the equations for \eqref{eq:vol2-divided-cumulants}.
The division is valid in the integral polynomial modules.
The linear component remains $P_x$; the unit assertion
follows from the corresponding assertion for cumulants.
\end{proof}

The quotients in \eqref{eq:vol2-divided-cumulants} are constructed
before reducing modulo any power $\hbar^M$. This order gives
maps at every finite quantum precision. Division by
$\hbar^{n-1}$ would not be well defined on an arbitrary
class in the truncated ring alone.

\subsection{The determinant force}

For units $a,b\in R_0$, take
\[
 S(x,y)=\tfrac12(a+gy)x^2+\tfrac12by^2.
\]
The conditional Gaussian formulas are
\begin{equation}\label{eq:vol2-determinant-force}
 \begin{aligned}
 Z_x&=(1+gy/a)^{-1/2},\\
 S_{\mathrm{eff}}&=\tfrac12by^2
                         +\tfrac{\hbar}{2}\log(1+gy/a),\\
 P_x(x^2)&=\frac{\hbar}{a+gy}.
 \end{aligned}
\end{equation}
All inverses and the square root are their $g$-adic expansions
with the displayed constant terms. The logarithmic determinant
changes the retained force by $\hbar g/(2(a+gy))$.
It is visible in a single cochain equation: on $\eta_yx^2$,
both $P_xD_S$ and $D_{S_{\mathrm{eff}}}P_x$ have value
\begin{equation}\label{eq:vol2-determinant-force-value}
 \frac{by\hbar}{a+gy}
                  +\frac{3g\hbar^2}{2(a+gy)^2}.
\end{equation}
On the source side use $P_x(x^4)=3\hbar^2/(a+gy)^2$.
On the target side differentiate $\hbar/(a+gy)$ and include
the determinant force. Omitting that force leaves the
nonzero defect $g\hbar^2/(2(a+gy)^2)$.

\subsection{A retained odd generator also changes}

The projection is accompanied by the corrected inclusion
$i'=(1+H_0\delta)^{-1}i_0$ from
\eqref{eq:vol2-perturbation-formulas}. For the formal potential
\[
 S(x,y)=\tfrac12ax^2+\tfrac12by^2+g\,\frac{yx^3}{6},
\]
this correction is already nonzero on $\eta_y$:
\begin{equation}\label{eq:vol2-cubic-antifield-lift}
 \begin{split}
 i'(\eta_y)={}&\eta_y
  -g\eta_x\left(\frac{x^2}{6a}+\frac{\hbar}{3a^2}\right)\\
 &+g^2y\eta_x\left(\frac{x^3}{12a^2}
                          +\frac{5\hbar x}{12a^3}\right)+O(g^3).
 \end{split}
\end{equation}
Indeed $\delta\eta_y=gx^3/6$, and the Gaussian homotopy
sends $x^3$ to $\eta_x(x^2/a+2\hbar/a^2)$.
Applying $H_0\delta$ a second time gives the displayed
quadratic-coupling term.

The same model has
\begin{align}
 Z_x&=1+\frac{5g^2\hbar y^2}{24a^3}
             +\frac{385g^4\hbar^2y^4}{1152a^6}+O(g^6),
             \label{eq:vol2-cubic-partition}\\
 S_{\mathrm{eff}}
 &=\tfrac12by^2-\frac{5g^2\hbar^2y^2}{24a^3}
                       -\frac{5g^4\hbar^3y^4}{16a^6}+O(g^6).
                       \notag
\end{align}
The Gaussian sixth and twelfth moments give $Z_x$.
Subtracting half the square of its order-two term in
$\log Z_x$ gives the fourth-order effective interaction.
Thus the state, the retained force and the lift of an odd
observable change together under elimination.

\section{Sewing the state}
\label{sec:vol2-state-sewing}

Conditional elimination provides maps on observables and an
effective weight for the next integration. Both enter sewing.
The quadratic calculation separates the covariance that crosses
the cut from the determinant that normalizes the result.

\subsection{Covariance and determinant}

Order the coordinates as retained $y$ followed by eliminated
$x$, and let
\[
 A=\begin{pmatrix}B&L\\L^{\mathrm t}&D\end{pmatrix},
 \qquad S=B-LD^{-1}L^{\mathrm t}.
\]
Assume that $D$ and $S$ are invertible. The matrix $S$ is the
\emph{Schur complement} of the eliminated block. Completing
the square gives
\begin{equation}\label{eq:vol2-schur-square}
 \begin{split}
 \tfrac12(y,x)A(y,x)^{\mathrm t}
   =\tfrac12y^{\mathrm t}Sy
    +\tfrac12(x+D^{-1}L^{\mathrm t}y)^{\mathrm t}
                         D(x+D^{-1}L^{\mathrm t}y).
 \end{split}
\end{equation}
Put
\[
 E=\begin{pmatrix}1\\-D^{-1}L^{\mathrm t}\end{pmatrix},
 \qquad
 C_0=\begin{pmatrix}0&0\\0&D^{-1}\end{pmatrix}.
\]
The map $E$ extends a retained configuration to the critical
point in its eliminated fibre.

\begin{proposition}\label{prop:vol2-schur-sewing}
One has
\begin{equation}\label{eq:vol2-schur-covariance}
 A^{-1}=C_0+ES^{-1}E^{\mathrm t},\qquad
 \det A=\det D\det S.
\end{equation}
The normalized Wick functional of covariance $\hbar A^{-1}$
equals successive Wick integration in
$\xi=x+D^{-1}L^{\mathrm t}y$ and then in $y$, with
covariances $\hbar D^{-1}$ and $\hbar S^{-1}$.
Successive Schur complements agree with elimination of the
combined block whenever the intermediate complements are
invertible.
\end{proposition}
\begin{proof}
The triangular matrix
\[
 M=\begin{pmatrix}1&0\\-D^{-1}L^{\mathrm t}&1\end{pmatrix}
\]
has determinant one and satisfies
$M^{\mathrm t}AM=\operatorname{diag}(S,D)$.
Taking inverses gives the first identity in
\eqref{eq:vol2-schur-covariance}; taking determinants gives
the second. The exponential generating function
\eqref{eq:vol2-wick-generating} for two independent Gaussian
blocks then becomes the generating function of covariance
$\hbar(C_0+ES^{-1}E^{\mathrm t})$. Equality of its polynomial
coefficients proves the state identity.

To compare successive block eliminations, solve the linear
critical-point equations in the eliminated variables. Solving
them one block at a time and solving the combined system give
the same solution. Substitution in the quadratic form therefore
gives the same retained form. The determinant identity and
the generating-function identity then compose in the same
order. This proves the assertion for any finite number of
admissible blocks.
\end{proof}

The change of odd generators is also fixed by this calculation.
For the coordinates $(y,\xi)$, write the odd generators as
$\vartheta_y,\vartheta_\xi$. The old ones become
\[
 \eta_{x_a}\longmapsto\vartheta_{\xi_a},\qquad
 \eta_{y_i}\longmapsto\vartheta_{y_i}
       +\sum_a(D^{-1}L^{\mathrm t})_{ai}\vartheta_{\xi_a}.
\]
These are the inverse-Jacobian transformations of coordinate
vector fields. Substitution in the Ward differential gives
the block-diagonal differential for
\eqref{eq:vol2-schur-square}. Thus the conditional comparison
also applies to all odd observables.

At a positive real specialization, the unnormalized generating
integral is
\begin{equation}\label{eq:vol2-unnormalized-gaussian}
 \begin{split}
 Z_A(J)
 &=\int_{\mathbb R^{r+s}}
        e^{(-q^{\mathrm t}Aq/2+J^{\mathrm t}q)/\hbar}\,dq\\
 &=(2\pi\hbar)^{(r+s)/2}(\det A)^{-1/2}
                       e^{J^{\mathrm t}A^{-1}J/(2\hbar)}.
 \end{split}
\end{equation}
The two identities in \eqref{eq:vol2-schur-covariance} preserve
respectively its moments and its normalization. Over a general
coefficient ring, an unnormalized state requires a specified
square-root normalization in addition to the algebraic
determinant identity. The normalized Wick functional itself
needs no such choice.

\subsection{The intermediate weight}

Return to the block potential
\eqref{eq:vol2-conditional-potential}. Let $P_y^0$ denote the
free Wick expectation in the next block. Its effective
conditional expectation is
\[
 P_y^{\mathrm{eff}}(b)
      =\frac{P_y^0(Z_xb)}{P_y^0(Z_x)}.
\]
Thus the factor $Z_x$, which may depend on $y$, is part of
the next integration.

\begin{theorem}\label{thm:vol2-integral-sewing}
Successive normalized elimination equals joint elimination:
\begin{equation}\label{eq:vol2-normalized-fubini}
 P_y^{\mathrm{eff}}P_x(f)
  =\frac{P_y^0P_x^0(fe^{-V/\hbar})}
         {P_y^0P_x^0(e^{-V/\hbar})}.
\end{equation}
For these maps, and for further admissible blocks, the divided
cumulants of \eqref{eq:vol2-divided-cumulants} compose by
\begin{equation}\label{eq:vol2-divided-sewing}
 \Phi_n^{QP}(b_1,\ldots,b_n)
   =\sum_{\pi\in\Pi_n}\epsilon(\pi;b)
      \Phi_{|\pi|}^Q
       \left(\Phi_{|E|}^P(b_j:j\in E):E\in\pi\right).
\end{equation}
These equalities hold over the integral quantum ring, and
iterated sewing is associative in every arity.
\end{theorem}
\begin{proof}
Insert the normalized numerator of $P_x$ into the next
expectation. Its denominator $Z_x$ cancels the intermediate
weight, giving \eqref{eq:vol2-normalized-fubini}.
Free Wick operations on separate blocks commute. At each
coupling order only finitely many polynomial contractions
occur, so this cancellation and Gaussian Fubini are defined
coefficientwise. Theorem~\ref{thm:vol2-integral-conditional}
shows that the resulting normalized maps are integral.

Equation~\eqref{eq:vol2-cumulant-sewing} gives the partition
composition law before division. In the term indexed by
$\pi$, the quantum exponents add to
\[
 (|\pi|-1)+\sum_{E\in\pi}(|E|-1)=n-1.
\]
Dividing that common power gives
\eqref{eq:vol2-divided-sewing}. Composition of the underlying
coalgebra maps is associative; the divided formula preserves
that equality. This argument is also the composition mechanism
for derived BV morphisms in Bandiera
\cite{Bandiera2020}*{\S2.2, Proposition 2.8}.
\end{proof}

For example, let $X=Y+\xi$, where the first expectation
eliminates a centred Gaussian $\xi$ of variance $\sigma$
and leaves $Y$. Its generating function on the square is
\[
 P(e^{tX^2})=(1-2\sigma t)^{-1/2}
                    \exp\!\left(\frac{tY^2}{1-2\sigma t}\right).
\]
The logarithm gives, in every arity,
\begin{equation}\label{eq:vol2-conditional-square-cumulants}
 \kappa_n^P(X^2,\ldots,X^2)
  =2^{n-1}(n-1)!
                    \bigl(\sigma^n+nY^2\sigma^{n-1}\bigr).
\end{equation}
If the next expectation eliminates an independent centred
Gaussian $Y$ of variance $\tau$, the joint variable $X$ has
variance $\sigma+\tau$. Equation~\eqref{eq:vol2-cumulant-sewing}
therefore combines all the conditional terms into
$2^{n-1}(n-1)!(\sigma+\tau)^n$.
Keeping only the conditional mean of $X^2$ would lose these
higher terms.

\section{Coordinate descent of the whole Ward complex}
\label{sec:vol2-ward-coordinates}

A change of coordinates transforms the integration density
and the odd observables at the same time. The logarithm of
the Jacobian enters the potential, while the inverse Jacobian
transforms the odd generators. Their two contributions are
required by the cochain equation.

Let $S(x)=\tfrac12x^{\mathrm t}Ax+V(x)$ be a finite potential
as above. Take a formal change of variables
$x=\phi(q)=q+O(g)$, polynomial in $q$ at each coupling order.
Its Jacobian $J=D_q\phi$ is invertible $g$-adically. Define
\begin{equation}\label{eq:vol2-coordinate-pullback}
 \Phi_\phi(f)=f\circ\phi,\qquad
 \Phi_\phi(\eta_i)=\sum_a(J^{-1})_{ai}\vartheta_a,
 \qquad
 S_\phi=S\circ\phi-\hbar\log\det J.
\end{equation}
Here $\vartheta_a$ are the new degree-minus-one generators.
The logarithm is the unique formal logarithm with zero
constant coupling term, since $\det J=1+O(g)$.

\begin{theorem}\label{thm:vol2-ward-coordinate-descent}
The map $\Phi_\phi$ is a unital graded-algebra and cochain
isomorphism over $R_0[\hbar][[g]]$, and it preserves the
normalized formal state:
\begin{equation}\label{eq:vol2-coordinate-state}
 D_{S_\phi}\Phi_\phi=\Phi_\phi D_S,
 \qquad p_{S_\phi}\Phi_\phi=p_S.
\end{equation}
For a second formal change $q=\psi(r)$,
\[
 (S_\phi)_\psi=S_{\phi\circ\psi},\qquad
 \Phi_{\phi\circ\psi}=\Phi_\psi\Phi_\phi.
\]
\end{theorem}
\begin{proof}
Regard the odd generators as coordinate vector fields and
their products as polyvectors. For the weighted volume form
$\rho=e^{-S/\hbar}dx_1\cdots dx_m$, divergence on polyvectors
has the coordinate formula
\[
 \operatorname{Div}_\rho
    =\sum_i\left(\partial_{x_i}-\hbar^{-1}S_i\right)
                                                    \partial_{\eta_i}.
\]
Thus $D_S=-\hbar\operatorname{Div}_\rho$.
This notation abbreviates the polynomial Ward operator; the
displayed product has no negative quantum power.
Contraction with a volume form identifies this divergence
with the exterior differential, with the exterior signs fixed
by left odd differentiation. It is therefore natural under
coordinate changes. In the present coordinates,
\[
 \phi^*\rho=e^{-S(\phi(q))/\hbar}\det J\,dq
                 =e^{-S_\phi(q)/\hbar}dq,
\]
and vector fields pull back by $J^{-1}$. These two facts
give the first equation in \eqref{eq:vol2-coordinate-state}
in every exterior degree.

For the state, interpolate by
$\phi_u(q)=q+u(\phi(q)-q)$, $0\leq u\leq1$.
The derivative of its pulled-back density times any pulled-back
observable is a total derivative in the integration variables.
At each coupling order, it is a polynomial times the free
Gaussian weight. Formal Gaussian integration by parts
annihilates that derivative. Both the unnormalized numerator
and denominator are consequently constant in $u$.
Their quotient proves the state equation on even observables;
negative cohomological degrees vanish on both sides.
The argument is formal in $g$ and uses no global real
invertibility of $\phi$.

Finally the chain rule gives
\[
 \det D(\phi\circ\psi)
       =((\det D\phi)\circ\psi)\det D\psi.
\]
The factors are commuting units with constant term one,
so their formal logarithms add. The inverse-Jacobian rule
has the same composition law on vector fields. This proves
the two cocycle identities.
\end{proof}

For $S=ax^2/2$ and $\phi(q)=q+gq^2$,
\begin{equation}\label{eq:vol2-nonlinear-jacobian}
 S_\phi=\tfrac a2(q+gq^2)^2-\hbar\log(1+2gq),
 \qquad
 \Phi_\phi(\eta)=\frac{\vartheta}{1+2gq}.
\end{equation}
If the logarithmic term is omitted, the cochain defect on
this odd observable is $2g\hbar/(1+2gq)^2$.
The inverse coordinate begins
$q=x-gx^2+2g^2x^3-5g^3x^4+\cdots$; its fifth-power term
is $14g^4x^5$, which has zero Gaussian expectation.
Consequently
\[
 p_{S_\phi}(q)=-\frac{g\hbar}{a}
                 -\frac{15g^3\hbar^2}{a^2}+O(g^5),
 \qquad p_{S_\phi}(\phi(q))=0.
\]
Without the Jacobian in the density, the last expectation
instead begins with $-2g\hbar/a$. The same missing term
is detected by a cochain and by an ordinary observable.

\subsection{A flat family of state operations}

Let $\phi_s$ depend on finitely many parameters $s_i$ with
commuting derivatives $\partial_i$. Its coefficients may be
formal, or smooth on a specified parameter domain; all
identities remain coefficientwise in $g$. Put
\[
 S_s=S_{\phi_s},\qquad
 w_i=(D_q\phi_s)^{-1}\partial_i\phi_s,
 \qquad \Psi_i=\sum_a w_i^a\vartheta_a.
\]
For a vector field $w$, its Lie derivative on polyvectors
acts by $\mathcal L_w f=w(f)$ on functions and
$\mathcal L_w\vartheta_a=-\sum_b(\partial_{q_a}w^b)\vartheta_b$
on coordinate vector fields. Extend it as a derivation.
Define the connection
\[
 \nabla_i=\partial_i-\mathcal L_{w_i}.
\]

\begin{proposition}\label{prop:vol2-flat-state-family}
The coordinate family satisfies
\begin{equation}\label{eq:vol2-flat-state-family}
 \begin{gathered}
 \partial_iS_s=D_{S_s}\Psi_i
                  =w_i(S_s)-\hbar\operatorname{div}w_i,\\
 [\nabla_i,D_{S_s}]=0,
 \qquad [\nabla_i,\nabla_j]=0,
 \qquad \partial_i p_s(F)=p_s(\nabla_iF).
 \end{gathered}
\end{equation}
Every cumulant is horizontal:
\begin{equation}\label{eq:vol2-horizontal-cumulants}
 \partial_i\kappa_n^{p_s}(F_1,\ldots,F_n)
  =\sum_j\kappa_n^{p_s}
       (F_1,\ldots,\nabla_iF_j,\ldots,F_n).
\end{equation}
The same statement holds for the integral divided cumulants.
\end{proposition}
\begin{proof}
Differentiation of pullback gives
$\partial_i\Phi_s=\mathcal L_{w_i}\Phi_s$.
Differentiating the cochain equation in
Theorem~\ref{thm:vol2-ward-coordinate-descent} proves
$[\nabla_i,D_{S_s}]=0$.
Commuting two parameter derivatives of $\Phi_s$ gives
\[
 \partial_iw_j-\partial_jw_i-[w_i,w_j]=0,
\]
which is the curvature equation for $\nabla$.
Differentiating the pulled-back density gives the first
identity in \eqref{eq:vol2-flat-state-family}.

On a scalar polynomial $F$, the Ward equation gives
\[
 p_s(FD_{S_s}\Psi_i)=\hbar p_s(w_iF),
 \qquad p_s(D_{S_s}\Psi_i)=0.
\]
Indeed $D_{S_s}(F\Psi_i)=F D_{S_s}\Psi_i-\hbar w_iF$,
and its expectation is zero. The derivative of the
normalized state is the covariance with $\partial_iS_s$
divided by $-\hbar$, together with $p_s(\partial_iF)$.
Substituting these Ward identities gives
$p_s(\partial_iF-w_iF)$. The expression is integral;
both sides vanish on negative cohomological degrees.

Apply this state equation to every block of
\eqref{eq:vol2-state-cumulant}. Since $\mathcal L_{w_i}$
is a derivation, each input is differentiated once with
the same partition sign. This proves
\eqref{eq:vol2-horizontal-cumulants}. Dividing the already
integral cumulants by their known quantum powers gives
the final assertion.
\end{proof}

For example,
\[
 \phi_{u,v}(q)=q+ugq^2+vgq^3,
 \qquad
 w_u=\frac{gq^2}{1+2ugq+3vgq^2},\quad
 w_v=\frac{gq^3}{1+2ugq+3vgq^2}
\]
satisfies the displayed curvature equation as a rational
identity, before its $g$-adic expansion. On parameter forms
of degree $p$, the total differential is
$d_\nabla+(-1)^pD_{S_s}$. Flatness and the commutator
equation make its square zero. Restriction to parameter
faces preserves the cochain and cumulant equations.

\subsection{Coordinate changes within an eliminated fibre}

For a conditional reduction, take
$(x,y)=(\phi(q,y),y)$ and write
$J_x=D_q\phi$, $B=D_y\phi$.
The full inverse Jacobian transforms the old odd generators as
\begin{equation}\label{eq:vol2-fibre-antifields}
 \begin{aligned}
 \eta_{x_i}&\longmapsto
                     \sum_a(J_x^{-1})_{ai}\vartheta_{q_a},\\
 \eta_{y_j}&\longmapsto\vartheta_{y_j}
                     -\sum_a(J_x^{-1}B)_{aj}\vartheta_{q_a}.
 \end{aligned}
\end{equation}
The density changes by $\det J_x$.

\begin{proposition}\label{prop:vol2-conditional-coordinate-square}
The fibre substitution preserves the conditional partition
factor $Z_x(y)$ and effective potential. If $P_q$ is the
transformed conditional Ward map, then
\[
 P_q\Phi_\phi=P_x
\]
on the entire complex. The same square commutes on every
cumulant and divided cumulant, and with subsequent eliminations.
\end{proposition}
\begin{proof}
Apply the formal change-of-variables argument of
Theorem~\ref{thm:vol2-ward-coordinate-descent}, with $y$
held as coefficients. It preserves the conditional numerator
and its normalization, hence also $S_{\mathrm{eff}}$.
Projection kills every term containing a $\vartheta_q$.
In the second line of \eqref{eq:vol2-fibre-antifields}, only
$\vartheta_{y_j}$ survives; the same holds term by term in
products of retained odd generators. This proves the cochain
square in every exterior degree. Substitution in the
partition formula proves its cumulant versions, and
Theorem~\ref{thm:vol2-integral-sewing} proves compatibility
with successive blocks.
\end{proof}

\section{Canonical families and their marked homotopies}
\label{sec:vol2-gauge-state-homotopies}

Coordinate descent transports a state by an algebra isomorphism.
A canonical variation of a BV weight gives a more general
comparison: the normalized elimination commutes with parallel
transport up to a specified cochain homotopy. Its higher
components have one additional marked insertion.

\subsection{The integration domain and the master weight}

Let $A,B$ be graded-commutative coefficient algebras over a
characteristic-zero field with central parameter $\hbar$.
Initially allow Laurent coefficients in $\hbar$; integral
restrictions will be specified for the finite Gaussian sector.
Let $\Delta_A,\Delta_B$ have degree one, square zero,
differential-operator order at most two, and value zero on
the unit. In the potential model our convention is
\[
 \Delta_A=-\sum_i\partial_{q_i}\partial_{\eta_i}.
\]
Choose a degree-zero linear integration map $\mathcal P$
on a specified domain of weights in extensions of $A,B$, satisfying
the Stokes equation
\begin{equation}\label{eq:vol2-bv-stokes}
 \Delta_B\mathcal P=\mathcal P\Delta_A.
\end{equation}
Multiplication by the weights and division by the effective
weights below are required to be defined on this domain.

In the finite potential model, the domain consists of Gaussian
weights times polynomial observables, expanded coefficientwise
in the positive coupling filtration. Integration means formal
Wick integration in the eliminated variables. Gaussian
integration by parts proves \eqref{eq:vol2-bv-stokes} there.
For an analytic integral the integration cycle, decay and
absence of boundary terms are hypotheses of the same equation.

A weight $U_A$ of cohomological degree zero satisfying
$\Delta_AU_A=0$ is a \emph{master weight}. Put
$U_B=\mathcal P U_A$ and assume that $U_A$ and $U_B$
are invertible on the chosen domain. Define
\begin{equation}\label{eq:vol2-normalized-bv-map}
 D_A=\hbar U_A^{-1}\Delta_AU_A,
 \qquad D_B=\hbar U_B^{-1}\Delta_BU_B,
 \qquad p(f)=U_B^{-1}\mathcal P(U_Af).
\end{equation}
Conjugation in the two differential formulas denotes
operators: for example $D_Af=\hbar U_A^{-1}\Delta_A(U_Af)$.
For $U_A=e^{-S/\hbar}$ with $S$ independent of odd variables,
this is exactly the Ward differential
\eqref{eq:vol2-potential-ward}.

\begin{proposition}\label{prop:vol2-normalized-bv-map}
The effective weight is a master weight, and $p$ is a unital
cochain map. Successive normalized integrations compose when
each integration uses the effective weight from the preceding
one.
\end{proposition}
\begin{proof}
Stokes gives $\Delta_BU_B=0$. Conjugation therefore gives
$D_A1=D_B1=0$ and $D_A^2=D_B^2=0$. Moreover
\[
 D_Bp
   =\hbar U_B^{-1}\Delta_B\mathcal P U_A
   =\hbar U_B^{-1}\mathcal P\Delta_AU_A=pD_A,
\]
and $p1=1$ by definition. If $U_1=\mathcal P_1U_A$ and
$U_2=\mathcal P_2U_1$, the composite normalized map is
\[
 U_2^{-1}\mathcal P_2
        \bigl(U_1U_1^{-1}\mathcal P_1(U_Af)\bigr)
       =U_2^{-1}\mathcal P_2\mathcal P_1(U_Af).
\]
This is normalization of the composite integration, with
the entire intermediate weight retained.
\end{proof}

Theorem~\ref{thm:vol2-integral-conditional} supplies the
quasi-isomorphism and integral formulas in the finite
potential sector. Proposition~\ref{prop:vol2-normalized-bv-map}
uses only Stokes and the stated weight domain; those
hypotheses alone assert a cochain map.

\subsection{Transporting a canonical variation}

Let $t$ be a parameter and $R_t\in A$ have degree $-1$.
Suppose
\begin{equation}\label{eq:vol2-canonical-master-family}
 \partial_tU_A=\Delta_A(R_tU_A),\qquad \Delta_AU_A=0.
\end{equation}
The integration map and the two BV operators remain fixed.
A varying cycle can be treated in this form only after
its pullback to a fixed domain has been specified.
Write $m_R$ for multiplication by $R$ and put
\[
 L_R=[\Delta_A,m_R]=\Delta_Am_R+m_R\Delta_A,
 \qquad r_t=p_t(R_t).
\]
By Stokes, the effective family obeys
\begin{equation}\label{eq:vol2-effective-master-family}
 \partial_tU_B=\Delta_B(r_tU_B).
\end{equation}

Set $w_A=U_A^{-1}\partial_tU_A$ and
$w_B=U_B^{-1}\partial_tU_B$. Define the degree-zero transports
\begin{equation}\label{eq:vol2-normalized-transports}
 v_A=U_A^{-1}L_RU_A-m_{w_A},\qquad
 v_B=U_B^{-1}[\Delta_B,m_r]U_B-m_{w_B}.
\end{equation}
A commutator of a second-order operator with multiplication
has order at most one. Its value on the unit here is $w$.
Subtracting that value therefore makes each $v$ a derivation.
Since $[\Delta,L_R]=0$, differentiation gives
\begin{equation}\label{eq:vol2-moving-differentials}
 \partial_tD_A=[v_A,D_A],\qquad
 \partial_tD_B=[v_B,D_B].
\end{equation}
Thus the connections $\partial_t-v_A$ and $\partial_t-v_B$
commute with their respective differentials.

\begin{theorem}\label{thm:vol2-canonical-state-homotopy}
The degree-minus-one map
\begin{equation}\label{eq:vol2-marked-linear-homotopy}
 h_t(f)=\frac{p_t(R_tf)-p_t(R_t)p_t(f)}{\hbar}
\end{equation}
satisfies
\begin{equation}\label{eq:vol2-state-gauge-equation}
 \partial_tp_t+p_tv_A-v_Bp_t=D_Bh_t+h_tD_A,
 \qquad h_t(1)=0.
\end{equation}
The formula holds on the declared integration domain.
For polynomial marked inputs in a formal Gaussian
presentation with no negative quantum powers, the quotient
in \eqref{eq:vol2-marked-linear-homotopy} is integral.
\end{theorem}
\begin{proof}
On unnormalized weights define the odd operator
$K=\mathcal Pm_R-m_r\mathcal P$.
Equation~\eqref{eq:vol2-bv-stokes} gives
\begin{equation}\label{eq:vol2-unnormalized-gauge-equation}
 \mathcal PL_R-L_r\mathcal P
                      =\Delta_BK+K\Delta_A.
\end{equation}
Differentiate $p_t=U_B^{-1}\mathcal P U_A$ and add
$p_tv_A-v_Bp_t$. The two logarithmic-derivative
multiplication terms cancel, leaving
\[
 U_B^{-1}(\mathcal PL_R-L_r\mathcal P)U_A.
\]
Conjugating \eqref{eq:vol2-unnormalized-gauge-equation}
by the weights and using
$h_t=\hbar^{-1}U_B^{-1}KU_A$ proves
\eqref{eq:vol2-state-gauge-equation}.
Normalization gives $K(U_A)=0$, hence $h_t1=0$.
In the formal Gaussian sector, the numerator is the
two-input cumulant $\kappa_2^{p_t}(R_t,f)$.
Its connected graphs have at least one quantum power by
the argument of Theorem~\ref{thm:vol2-integral-conditional}.
\end{proof}

When parallel transports $T_A(t,0),T_B(t,0)$ for these
connections exist in the specified parameter domain,
\[
 \widehat p_t=T_B(t,0)^{-1}p_tT_A(t,0)
\]
has derivative homotopic to zero. Its homotopy is
$T_B(t,0)^{-1}h_tT_A(t,0)$.
Integrating this expression along a defined interval gives
an endpoint cochain homotopy. This may be an ordinary
analytic interval or coefficientwise integration in a
positive formal filtration; the evolution and endpoint
evaluation are part of the chosen domain.

\subsection{One marked input in every arity}

Apply the coalgebra construction of
\eqref{eq:vol2-cumulant-coalgebra} to $p_t$.
Let $S^c(p_t;h_t)$ be its unary tangent map: it applies
$h_t$ to exactly one symmetric factor and $p_t$ to every
other factor, with the graded tensor signs. It is a
degree-minus-one coderivation along $S^c(p_t)$.
Define
\[
 \mathcal H_t=\tau_B^{-1}S^c(p_t;h_t)\tau_A.
\]
For a derivation $v$, write $\widetilde v$ for its unary
coderivation. Derivations commute with multiplication
in the sense $\widetilde v\tau=\tau\widetilde v$.

\begin{theorem}\label{thm:vol2-marked-cumulant-homotopy}
The full operation comparison satisfies
\begin{equation}\label{eq:vol2-coalgebra-gauge-equation}
 \partial_t\mathscr K(p_t)
  +\mathscr K(p_t)\widetilde v_A
  -\widetilde v_B\mathscr K(p_t)
     =\mathscr D_B\mathcal H_t+\mathcal H_t\mathscr D_A.
\end{equation}
Its homotopy has Taylor coefficients
\begin{equation}\label{eq:vol2-marked-cumulant-homotopy}
 (\mathcal H_t)_n(b_1,\ldots,b_n)
       =\hbar^{-1}\kappa_{n+1}^{p_t}(R_t,b_1,\ldots,b_n).
\end{equation}
The marked input $R_t$ is placed first; the partition
sign convention in \eqref{eq:vol2-state-cumulant}
determines all signs. In the formal Gaussian sector,
$(\mathcal H_t)_n$ is divisible by $\hbar^{n-1}$.
Dividing gives a homotopy for the divided state operations
without inverting $\hbar$.
\end{theorem}
\begin{proof}
Differentiation of $S^c(p_t)$ differentiates one factor
at a time. Equation~\eqref{eq:vol2-state-gauge-equation}
and the cochain equation for $p_t$ therefore give the
unary-coderivation version of
\eqref{eq:vol2-coalgebra-gauge-equation}.
Conjugate it by $\tau_A,\tau_B$, and use the derivation
identity for $v_A,v_B$. This proves the full equation.

To compute its Taylor coefficient, differentiate the
partition formula for $\kappa_n$ by replacing one factor
$p(\prod_{j\in E}b_j)$ by
\[
 h\left(\prod_{j\in E}b_j\right)
 =\hbar^{-1}\left[
       p\left(R\prod_{j\in E}b_j\right)
           -p(R)p\left(\prod_{j\in E}b_j\right)\right].
\]
The first terms enumerate partitions in which the block
of $R$ is not a singleton. For a partition of the original
inputs into $k$ blocks, the second term occurs $k$ times.
Its coefficient is
$-k(-1)^{k-1}(k-1)!=(-1)^k k!$, precisely the coefficient
of the added singleton block. Moving $R$ to its block
produces the Koszul sign prescribed by placing it first.
This proves \eqref{eq:vol2-marked-cumulant-homotopy}.

A connected graph with $n+1$ marked inputs, $v$ interaction
vertices and $E$ edges has $E\geq n+v$.
Its quantum power after the interaction factors is at
least $n$. Division by the explicit $\hbar$ in
\eqref{eq:vol2-marked-cumulant-homotopy} leaves
$\hbar^{n-1}$. The same arity count as in
Theorem~\ref{thm:vol2-integral-conditional} permits division
of the entire homotopy equation.
\end{proof}

These homotopies also sew. For successive eliminations
$p_t:A\to B$ and $q_t:B\to C$ with the same intermediate
transport, differentiation gives
\begin{equation}\label{eq:vol2-gauge-homotopy-sewing}
 h_t^{\mathrm{tot}}=q_th_t^p+h_t^qp_t,
 \qquad
 \mathcal H^{\mathrm{tot}}
    =\mathscr K(q)\mathcal H^p
                         +\mathcal H^q\mathscr K(p).
\end{equation}
For normalized integration of the same master family,
the intermediate generator is $r_t=p_t(R_t)$.
The two covariance expressions in the first formula then
cancel their intermediate $q_t(r_tp_t(f))$ term, leaving
the marked covariance for $q_tp_t$ itself.
Three eliminations give the same three homotopy terms
in either parenthesization. Thus the path comparison
and the state sewing have a common composition law.

\subsection{The total master equation on parameter faces}

Let $T$ be a finite-dimensional parameter simplex.
Place parameter forms to the left of coefficient
elements, and extend $\Delta_A$ by
\[
 \Delta_A(\alpha\otimes a)
            =(-1)^{|\alpha|}\alpha\otimes\Delta_Aa.
\]
Suppose a total-degree-zero weight has been supplied with
\begin{equation}\label{eq:vol2-total-master-weight}
 \mathcal U_A\in\Omega^\bullet(T)\widehat\otimes A,
 \qquad(d_T+\hbar\Delta_A)\mathcal U_A=0,
\end{equation}
and invertible zero-form weight on the integration domain.
Its pushforward $\mathcal U_B=\mathcal P\mathcal U_A$
satisfies the same total equation by Stokes. The full
weights are invertible wherever their zero-form parts
are, because positive form degree is nilpotent.

\begin{proposition}\label{prop:vol2-master-family-faces}
For a supplied family \eqref{eq:vol2-total-master-weight},
the operators
\[
 \mathbb D_A=\mathcal U_A^{-1}(d_T+\hbar\Delta_A)\mathcal U_A,
 \qquad
 \mathbb D_B=\mathcal U_B^{-1}(d_T+\hbar\Delta_B)\mathcal U_B
\]
and the normalized map
\[
 \mathbb p(F)=\mathcal U_B^{-1}\mathcal P(\mathcal U_AF)
\]
are compatible with restriction to every parameter face.
They give square-zero unital differentials and a cochain
map. Their cumulants supply the higher parameter
homotopies and their boundary equations.
\end{proposition}
\begin{proof}
The form differential and the extended BV operator
anticommute. Hence their sum squares to zero, and
conjugation preserves that identity. The total master
equation makes the conjugated differentials kill the unit.
Stokes gives $\mathbb D_B\mathbb p=\mathbb p\mathbb D_A$
by the calculation of
Proposition~\ref{prop:vol2-normalized-bv-map}.
All these operations commute with pullback to a face.
Apply the finite-length cumulant conjugation to this
cochain equation over the ground coefficient ring.
Its coefficients in parameter form degree are the
successive homotopy equations; integration of their
exterior derivatives gives their face boundaries by
Stokes on the parameter simplex.
\end{proof}

On an interval, \eqref{eq:vol2-canonical-master-family}
gives the actual total weight
\[
 \mathcal U_A=U_A+dt\,R_tU_A/\hbar.
\]
The coefficient of $dt$ in its normalized map is $h_t$,
and its total differential is
$D_A+dt(\partial_t-v_A)$ on zero-form inputs.
For two or more parameters, independently chosen
one-parameter generators need not extend to a total
master weight. If $\mathcal U_1$ is its proposed
one-form part, a two-form part must solve the degree-two
equation of \eqref{eq:vol2-total-master-weight}.
The existence of that part is an additional mathematical
condition. The flat coordinate families of
Proposition~\ref{prop:vol2-flat-state-family} give one
explicit class with all parameter compatibilities.

\subsection{A canonical shear with nonzero homotopy}

Take even variables $x,y$ and odd variables $\eta,\zeta$
of degree $-1$, with
\[
 \Delta=-\partial_x\partial_\eta-\partial_y\partial_\zeta.
\]
For invertible stiffnesses $a,b$, set
\begin{equation}\label{eq:vol2-shear-potential}
 S_t(x,y)=\tfrac12ax^2+\tfrac12b(y-tx)^2,
 \qquad d=a+bt^2,
 \qquad k=bt/d.
\end{equation}
Work formally near $t=0$, or with real $a,b>0$ and real
$t$. The odd generator $R=\zeta x$ satisfies
\[
 [\Delta,m_R]=-x\partial_y+\zeta\partial_\eta.
\]
Since $\partial_tS_t=-bx(y-tx)$, the weight
$U_A=e^{-S_t/\hbar}$ obeys
\eqref{eq:vol2-canonical-master-family}.

Complete the square in $x$, and normalize the eliminated
Gaussian relative to its value at $t=0$. This gives
\begin{equation}\label{eq:vol2-shear-state}
 \begin{aligned}
 p_t(x)&=ky,&
 \operatorname{Var}_{p_t}(x)&=\hbar/d,\\
 S_t^{\mathrm{eff}}(y)
   &=\frac{ab}{2d}y^2+\frac{\hbar}{2}\log(d/a),&
 r_t&=\zeta ky.
 \end{aligned}
\end{equation}
The effective differential and normalized transports are
\begin{equation}\label{eq:vol2-shear-transports}
 \begin{gathered}
 D_B=\left(\frac{ab}{d}y-\hbar\partial_y\right)\partial_\zeta,\\
 v_A=-x\partial_y+\zeta\partial_\eta,
 \qquad v_B=-ky\partial_y+k\zeta\partial_\zeta.
 \end{gathered}
\end{equation}
The determinant term is independent of $y$, so it does
not appear in $D_B$. It is nevertheless required by the
effective master-family equation
\eqref{eq:vol2-effective-master-family}: differentiating
the logarithmic determinant supplies its scalar term.

The map $p_t$ kills $\eta$ and retains $\zeta$.
For an even polynomial $f$, Gaussian integration by parts
gives
\begin{equation}\label{eq:vol2-shear-homotopy}
 h_t(f)=\frac{\zeta}{\hbar}
              \bigl(p_t(xf)-ky\,p_t(f)\bigr)
                  =\frac{\zeta}{d}p_t(\partial_xf).
\end{equation}
The homotopy is zero on terms containing either odd
generator. In particular $h_t(x)=\zeta/d$ is nonzero
and integral in $\hbar$.

On $x$, both sides of
\eqref{eq:vol2-state-gauge-equation} are
\[
 \partial_t(ky)-v_B(ky)=\frac{ab}{d^2}y
                       =D_B(\zeta/d).
\]
On the eliminated odd generator $\eta$, the left side
is $p_t(v_A\eta)=\zeta$, and the right side is
\[
 h_t(D_A\eta)=h_t(dx-bty)=\zeta.
\]
This equation sees a component that cannot be tested on
ordinary even observables alone.

The higher homotopy also acts on composites. If
$m=ky$ and $\sigma=\hbar/d$, then
\[
 \kappa_2(x^2,x^2)=4m^2\sigma+2\sigma^2,
 \qquad
 \kappa_3(x,x^2,x^2)=\sigma\partial_m\kappa_2(x^2,x^2)
                              =8m\sigma^2.
\]
The derivative identity follows by differentiating
the Gaussian generating function with respect to its mean.
Equation~\eqref{eq:vol2-marked-cumulant-homotopy} therefore
gives
\begin{equation}\label{eq:vol2-shear-higher-component}
 (\mathcal H_t)_2(x^2,x^2)
                         =\frac{8\zeta\,ky\,\hbar}{d^2}.
\end{equation}
The linear homotopy and this higher component are parts
of the same normalized comparison.

\section{Comparing normalized Gaussian choices}
\label{sec:vol2-gaussian-choices}

A fixed projection onto retained variables does not determine
Gaussian elimination. For linear observables $x,y$, the
covariance remains in
\[
 R(xy)=p(x)p(y)+\hbar P(x,y).
\]
Two choices with the same projection therefore require a
comparison on products. In a finite paired complex, a
contraction of the discarded coefficients constructs that
comparison and all its relative parameter extensions.

\subsection{A contractible complex of choices}

Let $(V,d)$ and $(H,d_H)$ be bounded complexes of finitely
generated projective $R_0$-modules. Fix a contraction
\[
 H\mathrel{\mathop{\rightleftarrows}^{i}_{p}}V,
 \qquad pi=1,\qquad dh+hd=1-ip,
 \qquad hi=ph=h^2=0.
\]
The maps $i,p$ are cochain maps and $h$ has degree $-1$.
Put $t=ip$ and $\mathsf E=t\otimes t$ on $V\otimes V$.
Every tensor of maps uses the graded rule
$(f\otimes g)(x\otimes y)=(-1)^{|g||x|}f(x)\otimes g(y)$.
For a homogeneous endomorphism $M$ of $V\otimes V$, define
\[
 \delta M=d_{V\otimes V}M-(-1)^{|M|}Md_{V\otimes V}.
\]

The symmetrized tensor contraction is
\begin{equation}\label{eq:vol2-reference-two-coefficients}
 A_0=\frac12\bigl(h\otimes1+t\otimes h
                         +1\otimes h+h\otimes t\bigr).
\end{equation}
The two unsymmetrized sums are exchanged by the graded
interchange of the tensor factors. Each has differential
$1-\mathsf E$, so
\begin{equation}\label{eq:vol2-normalized-two-coefficients}
 \delta A_0=1-\mathsf E,
 \qquad \mathsf EA_0=A_0\mathsf E=0.
\end{equation}
Let $\mathcal B$ be the graded endomorphisms commuting
with that interchange, with differential $\delta$, and set
\[
 \mathcal N=(1-\mathsf E)\mathcal B(1-\mathsf E).
\]
A \emph{normalized two-coefficient homotopy} is an element
$A\in\mathcal N^{-1}$ with $\delta A=1-\mathsf E$.
This condition imposes no equation on $A^2$.

\begin{lemma}\label{lem:vol2-gaussian-choice-contraction}
Left multiplication $S(M)=A_0M$ contracts
$(\mathcal N,\delta)$:
\[
 \delta S+S\delta=1.
\]
In particular, two normalized homotopies $A,B$ have the
specified comparison $L=A_0(B-A)$ with $\delta L=B-A$.
\end{lemma}
\begin{proof}
The graded Leibniz rule gives
\[
 \delta(A_0M)=(1-\mathsf E)M-A_0\delta M
                                      =M-A_0\delta M.
\]
Normalization is preserved: both
$\mathsf EA_0M$ and $A_0M\mathsf E$ vanish.
Composition also preserves commutation with the graded
interchange. Thus $S$ takes values in $\mathcal N$ and
the identity holds there. Apply it to the closed element
$B-A$ to obtain the last assertion.
\end{proof}

\subsection{The covariance homotopy}

Suppose a graded symmetric pairing
$b:V\otimes V\to R_0$ of degree $1$ is closed in the
Hom complex. Write $b_H=b(i\otimes i)$ for its retained
pairing. For any graded symmetric pairing $P$, let
$\partial_P$ denote unordered pair contraction on
$\operatorname{Sym}(V)$, normalized by
$\partial_P(xy)=P(x,y)$. On a product of homogeneous
linear observables, select two factors and move them
to the front with the Koszul sign. Apply $P$ to that pair
and retain the other factors in their order, then sum
over all selected pairs.

Extend $d$ to $\operatorname{Sym}(V)$ as a derivation.
The two paired differentials are
\[
 D_b=d+\hbar\partial_b,
 \qquad D_H=d_H+\hbar\partial_{b_H}.
\]
Closedness of $b$ gives $[d,\partial_b]=0$; constant
coefficient pair contractions supercommute, so
$\partial_b^2=0$. Thus both displayed differentials
square to zero. For a normalized homotopy $A$, put
\begin{equation}\label{eq:vol2-choice-gaussian-map}
 P_A=-bA,
 \qquad R_A=\operatorname{Sym}(p)e^{\hbar\partial_{P_A}}.
\end{equation}
The covariance has degree zero. Its exponential is a
finite sum on every polynomial: each contraction lowers
the number of linear factors by two. These maps are
defined over $R_0[\hbar]$ and extend by scalar extension to
$R_0[[\hbar]]$ without completing polynomial degree.

\begin{proposition}\label{prop:vol2-choice-gaussian-homotopy}
Each $R_A$ is a cochain map from $D_b$ to $D_H$.
For two normalized homotopies $A,B$, define
\[
 L=A_0(B-A),\qquad Q=bL,\qquad
 P_u=(1-u)P_A+uP_B.
\]
Then the degree-minus-one map
\begin{equation}\label{eq:vol2-choice-gaussian-homotopy}
 K_{B,A}=\hbar\operatorname{Sym}(p)\partial_Q
                         \int_0^1e^{\hbar\partial_{P_u}}\,du
\end{equation}
satisfies
\[
 D_HK_{B,A}+K_{B,A}D_b=R_B-R_A.
\]
It vanishes on constants and on polynomials whose
coefficients all lie in $i(H)$.
\end{proposition}
\begin{proof}
The Hom differential obeys the graded composition rule.
Since $|b|=1$ and $\delta b=0$,
\[
 \delta P_A=b(1-\mathsf E)=b-p^*b_H.
\]
Here $p^*b_H=b_H(p\otimes p)$.
The identity $[d,\partial_{P_A}]=\partial_{\delta P_A}$
then gives
\[
 e^{\hbar\partial_{P_A}}D_b
       =(d+\hbar\partial_{p^*b_H})
                                      e^{\hbar\partial_{P_A}}.
\]
There are no further commutator terms, because the
constant coefficient contractions commute with the
degree-zero operator $\partial_{P_A}$.
Application of $\operatorname{Sym}(p)$ proves the
cochain equation for $R_A$.

By Lemma~\ref{lem:vol2-gaussian-choice-contraction},
$\delta L=B-A$. The same graded composition rule gives
\[
 \delta Q=-b(B-A)=P_B-P_A.
\]
Place parameter forms before coefficient maps. The
total-degree-zero covariance $P_u+du\,Q$ therefore
satisfies
\[
 (d_u+\delta)(P_u+du\,Q)=b-p^*b_H.
\]
Its exponential gives a total cochain map by the preceding
calculation. Its coefficient of $du$ is the integrand
in \eqref{eq:vol2-choice-gaussian-homotopy}. The
coefficient cochain equation, integrated from $0$ to $1$,
is the claimed homotopy identity. This integration is
polynomial integration over the rational ground ring.
Finally $L\mathsf E=0$, so $Q$ vanishes on two retained
coefficients; pair contraction also kills constants.
\end{proof}

The factors left after contraction matter. A quadratic
input has homotopy $\hbar Q(x,y)$; a cubic input has terms
$\hbar Q(x,y)p(z)$ and the other two pairings, with their
Koszul signs. A five-dimensional complex makes both
values nonzero. Take a basis
\[
 |u|=-1,\quad |v|=|r|=|e|=0,\quad |s|=1,
 \qquad du=v,\quad dr=s,
\]
with all other differentials zero. Retain $e$, and set
$h(v)=u$, $h(s)=r$, with $h$ zero on the other basis
vectors. Let $b(r,u)=b(u,r)=1$ and let its other basis
values vanish. This is a closed degree-one pairing.

Define the degree-minus-two endomorphism $J$ of
$V\otimes V$ by
\[
 J(s\otimes v)=J(v\otimes s)
                  =\tfrac12(r\otimes u+u\otimes r),
\]
and zero on other tensor basis vectors. It is normalized
and commutes with graded interchange. Thus
$A=A_0$ and $B=A_0+\delta J$ are normalized homotopies.
Their covariance data are
\begin{equation}\label{eq:vol2-choice-five-dimensional}
 P_A(r,v)=-\tfrac12,
 \qquad P_B(r,v)=\tfrac12,
 \qquad Q(s,v)=1.
\end{equation}
Indeed \eqref{eq:vol2-reference-two-coefficients} gives
$A_0(r\otimes v)=r\otimes u/2$, while
$\delta J(r\otimes v)=-(r\otimes u+u\otimes r)/2$.
Also $A_0\delta J(s\otimes v)
=(r\otimes u+u\otimes r)/2$, proving the third value.
Since $D_b(rv)=sv$ and $b_H=0$, the homotopy equation reads
\[
 R_B(rv)-R_A(rv)=\hbar=K_{B,A}(sv).
\]
After multiplication by the retained coefficient,
$K_{B,A}(sve)=\hbar e$.

\subsection{Prescribed parameter faces}

For the cube $I^m$, use polynomial parameter forms
\[
 \Omega_{\mathrm{poly}}^\bullet(I^m)
       =\mathbb Q[t_1,\ldots,t_m,dt_1,\ldots,dt_m],
\]
where the $dt_i$ are exterior generators of degree $1$.
A parameter family of normalized homotopies is a
total-degree-minus-one element $\mathcal A$ with
coefficients in $\mathcal N$ satisfying
\begin{equation}\label{eq:vol2-choice-total-equation}
 (d_{\mathrm{par}}+\delta)\mathcal A=1-\mathsf E.
\end{equation}
At a vertex, it is an ordinary normalized homotopy.
Restriction to a face sets its fixed coordinates to
$0$ or $1$ and their differentials to zero.

\begin{theorem}\label{thm:vol2-choice-relative-fillers}
A family satisfying \eqref{eq:vol2-choice-total-equation}
on a union of codimension-one faces, with agreeing
restrictions on intersections, extends to the cube.
Two extensions admit such a family between them
relative to the prescribed faces. The same assertion
holds after adjoining any finite number of parameters.
\end{theorem}
\begin{proof}
First extend the forms without their differential
equation. For a face $t_i=\varepsilon$, use its coordinate
projection and the scalar factor $1-t_i$ if
$\varepsilon=0$, or $t_i$ if $\varepsilon=1$.
Sum these extensions and subtract repeated extensions
from intersections by inclusion--exclusion. Opposite
faces have disjoint restriction conditions, since the
factor for one vanishes on the other. On each prescribed
face, terms from every other compatible face cancel
against their intersection terms. The resulting
polynomial form $\widetilde{\mathcal A}$ has the given
restrictions and coefficients in $\mathcal N$.

Write $D=d_{\mathrm{par}}+\delta$ and set
\[
 Z=D\widetilde{\mathcal A}-(1-\mathsf E).
\]
This is closed and restricts to zero on the prescribed
faces. Extend the contraction of
Lemma~\ref{lem:vol2-gaussian-choice-contraction} by
\[
 S_{\mathrm{par}}(\alpha\otimes M)
             =(-1)^{|\alpha|}\alpha\otimes A_0M.
\]
The parameter differential anticommutes with this
degree-minus-one map, and the coefficient identity
gives $DS_{\mathrm{par}}+S_{\mathrm{par}}D=1$.
Consequently
\begin{equation}\label{eq:vol2-choice-relative-correction}
 \mathcal A=\widetilde{\mathcal A}-S_{\mathrm{par}}Z
\end{equation}
solves \eqref{eq:vol2-choice-total-equation} and keeps
all the prescribed face values. To compare two
extensions, prescribe them at the ends of one more
interval and keep their common boundary family
constant along that interval. Apply the same correction.
Repeating this procedure gives the higher comparisons.
\end{proof}

Each family gives the total Gaussian map
$\operatorname{Sym}(p)\exp(\hbar\partial_{-b\mathcal A})$.
Composition with $b$ includes its sign past parameter
forms. The covariance calculation in
Proposition~\ref{prop:vol2-choice-gaussian-homotopy}
proves its total cochain equation. Thus the construction
supplies relative homotopies whenever the prescribed
Gaussian comparisons have these normalized coefficient
lifts. An arbitrary prescribed homotopy of observables
need not come with such a lift.

\section{From elimination to supported field operations}
\label{sec:vol2-state-geometric-comparison}

The finite constructions retain three kinds of data:
the relative support class in the normal direction,
the normalized state and all its cumulants, and the
parameter homotopies transporting that state. Conditional
elimination and coordinate changes construct these data
over integral quantum coefficients. A supplied total
master weight extends them over parameter faces, while
the normalized Gaussian choices above provide an explicit
class for which every relative extension exists.

The corresponding field-theoretic comparison requires a
common supported operation complex containing the
regular, meromorphic and contact coefficients of the
previous chapters. Its normal reduction must intertwine
the differential and every allowed collision operation;
its state map must carry the local vacuum and the
prescribed sewing and parameter homotopies. A regulated
master family must make its boundary terms and parameter
face equations meaningful on that same coefficient
domain. The finite endomorphism contraction above does
not yet act on those distributional products. These are
the missing constructions for a state-preserving
geometric comparison.

There is a second source of higher information when
observables obey algebraic constraints. Before applying
elimination to matrix observables, one must retain the
relations among their constraint equations. The resulting
complex is the starting point for Hamiltonian actions on
derived matrix spaces.
